\documentclass[manuscript,screen,nonacm]{acmart}
\usepackage{booktabs}
\settopmatter{printacmref=false}
\renewcommand\footnotetextcopyrightpermission[1]{}
\newcommand{\repo}{https://github.com/keshavdubeyy/reproagent_autose}
\newcommand{\rlink}[2]{\href{\repo/blob/main/#1}{#2}}

\begin{document}

\title{ReproAgent: Reproducing and Fixing Python Bugs from Bug Reports with a Local Small Language Model}
\subtitle{CS6.403 Autonomous Software Engineering --- Mid-Submission Progress Report}

\author{Keshav Dubey}
\author{Rishabh Singh}
\author{Aasritha Kalluri}
\affiliation{\institution{IIIT Hyderabad}\country{India}}

\begin{abstract}
Before a bug can be fixed it must be reproduced, usually as a failing test, and this step is still largely manual.
ReproAgent is a coding agent that runs entirely on a local small language model (\texttt{qwen3:8b}~\cite{qwen3} via Ollama~\cite{ollama}).
It reads a natural-language bug report, explores the repository, writes a pytest test that fails because of the bug,
and, following reviewer feedback, then proposes a patch and verifies that it makes the test pass without breaking the existing suite.
Our central design choice is that the \emph{harness}, not the model, decides whether each stage succeeded:
only three steps of the pipeline use the model, while all checks are deterministic or human.
We report how we built it step by step, the architecture, a benchmark of eight real BugsInPy bugs plus a toy bug, a working chat interface,
and a pilot evaluation (13 of 27 planned runs). The full agent produces a harness-valid failing test for 7 of 9 bugs,
but only 3 of these pass the fail-to-pass oracle, which makes false reproductions our main open problem.
\end{abstract}

\maketitle

\noindent\textbf{Artifact.} Code, data, results and documentation: \url{\repo}.
Run \texttt{make setup \&\& make test \&\& make demo} to reproduce the toy example.

\section{Problem and Motivation}
Developers receive bug reports as prose. Turning that prose into an executable failing test is the first step of any fix,
and the test also guards against the bug coming back. Large hosted models can do this~\cite{libro}, but they are expensive, send code to third parties,
and consume significant energy. We ask whether a \emph{small, local} model can do it reliably when it is wrapped in a carefully engineered harness,
as encouraged by the course~\cite{course}.

\textbf{Scope change after review.} The proposal was accepted with the comment ``make the harness also fix the bug''.
We added a Fix stage that only counts as successful once the harness has verified it
(\rlink{docs/adr/0005-add-fix-stage.md}{ADR-0005}).

\textbf{Hypotheses.}
\begin{itemize}
\item \textbf{H1}: an agent with repository tools reproduces more bugs than generating a test from the description alone.
\item \textbf{H2}: an execution-feedback loop with retries beats a single generation.
\item \textbf{H3}: with harness-owned verification, a local small model can produce fixes that pass the developers' own tests.
\end{itemize}

\section{Related Work}
\textbf{Bug reproduction.} LIBRO~\cite{libro} showed that LLMs can turn bug reports into reproducing tests, using a large hosted model with
sampling and ranking. SWT-Bench~\cite{swtbench} frames test generation from issues as a benchmark and uses the fail-to-pass criterion we
adopt as our oracle.

\textbf{Repair agents.} SWE-bench~\cite{swebench} established issue resolution on real repositories, and SWE-agent~\cite{sweagent}
showed that a purpose-built agent--computer interface (bounded search, file viewer, guarded editor) matters as much as the model; our tools follow
that design. Both mostly assume large hosted models.

\textbf{Agent design.} We build on ReAct~\cite{react} for the loop, on the pattern catalogue
of Liu et al.~\cite{liu-patterns} for choosing which patterns to use or avoid, and on EARS~\cite{ears} for requirements.

\textbf{Context files.} Evidence that repository context files raise cost more than success~\cite{agentsmd} motivates keeping our prompts and \texttt{AGENTS.md} short.

\textbf{Our difference.} a \emph{local small model}, reproduction \emph{and} repair in one pipeline, and every success decision owned by the harness.

\textbf{Benchmark.} Benchmark bugs come from BugsInPy~\cite{bugsinpy}.

\section{How We Built It, Step by Step}
\textbf{Step 1 --- Understand the course and frame the problem.} We read the course outline, project guidelines and lectures, and wrote a proposal
for an agent that turns bug reports into failing tests on a local model. It was accepted with the request to also fix the bug.

\textbf{Step 2 --- Decide where AI is needed.} Following the lecture decomposition of autonomous bug fixing into bounded contexts,
we listed every part of the job and asked, for each, whether a rule could do it (Section~\ref{sec:plan}, Table~\ref{tab:ai}).
The answer gave three AI steps and moved every success decision into code.

\textbf{Step 3 --- Specify before coding.} We wrote requirements in EARS form (\rlink{docs/spec.md}{\texttt{spec.md}}, FR-01--FR-10),
a short \rlink{AGENTS.md}{\texttt{AGENTS.md}}, and decision records for the main choices: local model (ADR-0001), harness-owned verification (ADR-0002),
temp-copy sandbox (ADR-0003), tool design (ADR-0004) and the fix stage (ADR-0005), all in \rlink{docs/adr}{\texttt{docs/adr/}}.

\textbf{Step 4 --- Build a prototype on a toy bug.} We implemented the model client (\rlink{reproagent/llm.py}{\texttt{llm.py}}), the sandbox and tools
(\rlink{reproagent/workspace.py}{\texttt{workspace.py}}), the ReAct loop (\rlink{reproagent/agent.py}{\texttt{agent.py}}) and the two-stage pipeline
(\rlink{reproagent/pipeline.py}{\texttt{pipeline.py}}), and tested them on a seeded median bug (\rlink{examples}{\texttt{examples/}}).
The agent reproduced and fixed it in one attempt.

\textbf{Step 5 --- Add real bugs.} We built a benchmark from BugsInPy (\rlink{benchmarks}{\texttt{benchmarks/}}): buggy and fixed snapshots,
hand-written reports, and the developers' tests. This took longer than planned: projects need Python~3.8, per-bug (not per-project) requirements,
extra packages such as \texttt{nose} and \texttt{setuptools}, and old pins that do not build on Apple Silicon. A shared environment was polluted by one bug,
so we moved to one environment per bug and kept only bugs whose developer test is verifiably fail-to-pass.

\textbf{Step 6 --- Pilot and harden the harness.} A first pilot exposed harness bugs, each fixed with a guard:
pytest usage errors (exit code 4) were wrongly accepted as reproductions, so now only exit code 1 counts;
the agent declared ``needs more info'' too eagerly, so it became a soft signal followed by a retry;
the agent wrote tests asserting the \emph{buggy} behaviour, so the prompt forbids it and the oracle catches it;
edits failed on duplicated lines, so the editor takes a line number; and the fix stage once edited an existing test, so test files became read-only.

\textbf{Step 7 --- Build a usable product.} We added a CLI and a local chat UI modelled on the course lab's tool (\rlink{docs/product-flows.md}{product flows}):
connect a repository, describe the bug or paste a GitHub issue, confirm the agent's understanding, watch reproduction live, then fix and apply with approval.
Building it surfaced more guards: a busy lock against double clicks, an ``interrupted'' state after restarts, and a check that the repository
has not changed before applying.

\textbf{Step 8 --- Align with the lectures and document.} We mapped every course concept to the code (Section~\ref{sec:concepts},
\rlink{docs/design-overview.md}{design overview}), drew the architecture (\rlink{docs/architecture.md}{architecture}),
and wrote the status and this report. We paused the full evaluation to finish these first.

\section{How We Planned the Agent}\label{sec:plan}
We started from the job (report $\rightarrow$ failing test $\rightarrow$ patch $\rightarrow$ evidence), split it into bounded contexts
(issue management, bug resolution, verification), and for each part asked whether it needs AI at all.
The model is given only the fuzzy parts (understanding prose, deciding where to look, writing code);
everything that can be checked is checked by code (Table~\ref{tab:ai}). This follows the view that \emph{agent = model + harness}:
we never change the model, so all of our engineering is in the harness.

\begin{table}[t]
\caption{Pipeline parts and whether they use AI.}
\label{tab:ai}
\small
\begin{tabular}{@{}lll@{}}
\toprule
Part & What it does & Kind \\
\midrule
Issue retrieval & text, .md, or GitHub issue URL & tool \\
Bug specification & report $\rightarrow$ structured JSON & \textbf{AI} \\
Confirm understanding & user edits/accepts & human \\
Repository explorer & list, search ($\le$50 hits), read & tool \\
Test designer & explore + write repro test & \textbf{AI} (ReAct) \\
Test executor & pytest in sandbox, timeout & tool \\
Result analyzer & real failure vs.\ import/syntax error & deterministic \\
Fix generator & edit source (tests read-only) & \textbf{AI} (ReAct) \\
Verification & repro passes, no new failures & deterministic \\
Apply & write test+patch to user repo & human + tool \\
\bottomrule
\end{tabular}
\end{table}

\section{Architecture}
ReproAgent is about 1.1k lines of standard-library Python plus pytest, in four layers (diagrams in \rlink{docs/architecture.md}{\texttt{architecture.md}}).
\textbf{Interface}: a CLI (\texttt{run}, \texttt{eval}) and a local chat UI.
\textbf{Application}: the pipeline (stage order, prompts, retries, all pass/fail decisions) and a session state machine.
\textbf{Agent}: a generic ReAct loop that sends tool schemas to the model, executes the chosen tool, and returns the result
(errors are returned as observations rather than crashes), with a step budget and a timestamped trace.
\textbf{Execution}: a \texttt{Workspace} that copies the repository into a temporary sandbox and exposes path-confined tools
in the style of SWE-agent's agent--computer interface: bounded search, a line-numbered file viewer, an editor that requires a unique match
(or an explicit line number to disambiguate), and a pytest runner with a scrubbed environment.

\textbf{Nested loops.} Tool steps ($\le$15) run inside attempts ($\le$3 per stage, reverting the sandbox between attempts and feeding
the previous failure back as feedback), inside stages (Reproduce, Fix), inside human checkpoints (confirm understanding, apply patch).

\textbf{Verification.} A reproduction is accepted only if pytest exits with code 1 (a genuine assertion or runtime failure);
collection, import and syntax errors are rejected. A fix is accepted only if the reproduction test now passes and no test that passed
before the fix now fails. Test files are read-only during Fix, so the agent cannot ``fix'' the bug by editing tests.
We deliberately avoid model-as-judge: an 8B judge is less reliable than pytest.
For evaluation we add a \emph{fail-to-pass oracle}~\cite{swtbench}: the generated test must also pass on the developers' fixed version.

\textbf{Safety.} The user's repository is read-only until the human approves \emph{Apply}, which refuses if the repository changed
since the session started. The UI is bound to localhost and no code leaves the machine.

\section{Mapping to Course Concepts}\label{sec:concepts}
Table~\ref{tab:concepts} lists each lecture concept and how ReproAgent uses it, or why it does not.

\begin{table}[t]
\caption{Course concepts in ReproAgent (Y used, P planned, N deliberately not used).}
\label{tab:concepts}
\footnotesize
\begin{tabular}{@{}p{2.6cm}cp{4.6cm}@{}}
\toprule
Concept & & How / why \\
\midrule
ReAct, tool calling & Y & one loop reused for both stages \\
Agent = model + harness & Y & model frozen; all work in the harness \\
Agentic vs.\ fixed workflow & Y & fixed outer pipeline, agentic inside stages \\
Temperature, prompting & Y & temp 0.2, zero-shot, prompt chaining, JSON output \\
Chain-of-thought & N & reasoning tokens off: time and energy on an 8B model \\
Loop engineering & Y & tool/attempt/stage/human loops with budgets \\
Verification types & Y & deterministic, state, grounding, human \\
Model-as-judge & N & pytest is a more reliable judge \\
SWE-agent interface & Y & bounded search, viewer, guarded editor, sandbox \\
Reflection & Y & failed attempt's output and diff fed back \\
Human-in-the-loop & Y & confirm understanding, approve Apply \\
Proactive goal creator & Y & intake lists missing information \\
Voting, debate & N & multiply cost; tests already decide \\
Sub-agents & N & one task, one loop \\
Memory & Y & working (messages), episodic (feedback), sessions \\
RAG & N & live search beats a stale index for one repo \\
Specs, EARS, AGENTS.md & Y & \texttt{spec.md}, short \texttt{AGENTS.md} \\
ADRs & Y & five decision records \\
MCP & P & expose our tools as an MCP server \\
Skills & P & one \texttt{SKILL.md}, evaluated with/without \\
Eval: baseline delta, cost & Y & three modes; tokens and time logged \\
\bottomrule
\end{tabular}
\end{table}

\section{Evaluation Setup}
\textbf{Benchmark.} A toy median bug plus eight BugsInPy bugs~\cite{bugsinpy}: tqdm-1, 3, 4, 7, 9; thefuck-1, 2; httpie-5
(\rlink{benchmarks/dataset.json}{\texttt{dataset.json}}, reports in \rlink{benchmarks/reports}{\texttt{benchmarks/reports/}}).
\textbf{Modes.} \texttt{full} (agent loop with tools and retries), \texttt{single\_shot} (one generation from the report plus the file list),
\texttt{desc\_only} (report only). All use the same model and settings, so differences come from the harness alone.
\textbf{Metrics.} Harness-valid reproduction; true reproduction (fail-to-pass); fix accepted by the harness;
fix passes the developers' test; LLM calls, tokens and wall time per run (wall time and tokens drive energy).

\section{Pilot Results}
Table~\ref{tab:pilot} summarises the 13 runs completed so far (Apple Silicon laptop; raw JSON in \rlink{results/pilot}{\texttt{results/pilot/}}).

\begin{table}[t]
\caption{Pilot results (13 of 27 runs).}
\label{tab:pilot}
\small
\begin{tabular}{@{}lrrrrrr@{}}
\toprule
Mode & $n$ & Valid & F2P & Fixed & Dev test & Time (s) \\
\midrule
full & 9 & 7 & 3 & 2 & 1 & 383 \\
single\_shot & 4 & 2 & 2 & 1 & 0 & 369 \\
desc\_only & 0 & -- & -- & -- & -- & -- \\
\bottomrule
\end{tabular}
\end{table}

The full agent fixed the toy bug and tqdm-4; the tqdm-4 patch also passes the developers' test (9 LLM calls, 101~s).
On tqdm-1 and tqdm-4, \texttt{single\_shot} failed where \texttt{full} succeeded, consistent with H2, while on tqdm-3 the reverse held.
The single-shot fix for tqdm-3 passed our checks but failed the developers' test, an example of a patch overfitting to its own test.
The most important finding is that 4 of the 7 harness-valid reproductions (tqdm-7, tqdm-9, thefuck-2, httpie-5) \emph{fail the oracle}:
they fail on the buggy code but also on the fixed code, so they do not capture the bug. Our harness check alone cannot detect this,
and those runs then consumed the full fix budget (40--51 LLM calls, up to 14 minutes). The sample is too small for conclusions about H1--H3.

\section{Status}
\textbf{Done.} Reproduce and Fix pipeline with retries and harness checks; sandbox; three evaluation modes; nine-bug benchmark
with per-bug environments; CLI and chat UI verified end to end; documentation (spec, ADRs, architecture, design overview);
12 unit tests that need no model; public repository with contribution guidelines.

\textbf{Remaining.} Complete the 27 runs, with three seeds each for variance; grow the benchmark to 15--20 bugs;
a failure taxonomy for every unsuccessful run; reducing false reproductions; the MCP server and skill ablation;
the final report, demo and viva.

\section{Challenges and Next Steps}
\textbf{False reproductions} are the main problem: a small model often writes a test that fails for an unrelated reason
(wrong setup, asserting the buggy behaviour, or a mis-located symptom). We plan to require the test to check the report's
\emph{expected} behaviour explicitly, add a cheap self-check that compares the failure message with the reported symptom,
and stop the Fix stage early when the reproduction looks weak.
\textbf{Speed} limits scale: runs take 1.5--14 minutes, so a full evaluation takes hours; we will cap budgets on hopeless runs
and report cost next to quality.
\textbf{Legacy environments} for BugsInPy remain fragile and restrict which bugs we can use.
\textbf{Small-model tool use} is brittle (early \texttt{needs\_info}, repeated calls, ambiguous edits); each failure mode
we found became a harness guard, and we will quantify how often each guard fires.

\bibliographystyle{ACM-Reference-Format}
\bibliography{references}

\end{document}
